% ===================================================================
%  QUADRO N.º 12 — A Estação. — Os Caminhos de ferro. — Os Barcos.
%
%  Ceci n'est PAS une transcription : c'est une traduction, faite en
%  2026, en portugais d'aujourd'hui. Voir texto/pt/10-quadro-01.tex
%  pour la consigne, qui vaut pour les seize fichiers.
% ===================================================================

%%K t12-apar-1 apar
\VUsaut{4.0mm}
\VUtitre{15.0pt}{60}{QUADRO N.º 12}
\VUsaut{2.4mm}
\VUpk{11.4pt}{40}{A Estação. --- Os Caminhos de ferro. --- Os Barcos.}
\VUsaut{3.4mm}
\textsc{Nota.} --- \textit{Os horários usados em cada país são
diferentes, pelo que tomámos como exemplo as horas do} \VUgras{quadro
francês}.

%%K t12-01-1 p
1. --- Acabo de fazer uma viagem pela Alemanha. Antes de a empreender
tinha comprado um \VUgras{guia de horários} \textsuperscript{(15)} para saber a
hora de partida dos comboios. Tendo verificado que o comboio mais cómodo
saía de Paris às nove da noite e chegava a Estrasburgo à uma e treze,
preparei a minha viagem e, no dia seguinte, mandei levar-me à estação.

%%K t12-02-1 p
2. --- Quando entrei na grande sala das partidas, atrás de um velho
\VUgras{pároco} \textsuperscript{(2)} de aldeia que levava na mão a sua \VUgras{bolsa
de viagem} \textsuperscript{(3)}, já tinham chegado muitos \VUgras{viajantes}
\textsuperscript{(1)}.

%%K t12-02-2 p
Como queria mandar um \VUgras{telegrama} \textsuperscript{(5)}, pedi a um
\VUgras{revisor} \textsuperscript{(6)} que me indicasse a \VUgras{estação
telegráfica} \textsuperscript{(4)}.

%%K t12-03-1 p
3. --- Antes de passar à \VUgras{sala de espera} \textsuperscript{(7)} fui tirar um
\VUgras{bilhete} \textsuperscript{(9)} de ida e volta.

%%K t12-04-1 p
4. --- Não esperei muito na \VUgras{bilheteira} \textsuperscript{(8)}, porque havia
pouca gente. Um \VUgras{soldado} \textsuperscript{(10)} que saía de licença cedeu-me
amavelmente a sua vez, de modo que tive tempo de pedir alguns dados ao
\VUgras{chefe de estação} \textsuperscript{(13)}, que falava com o
\VUgras{comissário} \textsuperscript{(12)} de vigilância e com o \VUgras{gendarme}
\textsuperscript{(11)} de serviço.

%%K t12-tit-1 sub
\VUpk{10.2pt}{40}{Despachar a bagagem.}
\VUsaut{3.0mm}

%%K t12-05-1 p
5. --- Depois de comprar um jornal no \VUgras{quiosque de livros}
\textsuperscript{(14)}, mandei pesar a minha \VUgras{bagagem} \textsuperscript{(17)}, que
um \VUgras{moço} \textsuperscript{(16)} acabava de trazer num \VUgras{carrinho}
\textsuperscript{(19)}; o \VUgras{empregado} \textsuperscript{(22)} encarregado da
\VUgras{báscula} \textsuperscript{(21)} disse-me que a minha mala pesava trinta e
cinco quilos, o que fazia um \VUgras{excesso} de cinco quilos, pelo qual
tinha de pagar um suplemento no \VUgras{guichê das bagagens}
\textsuperscript{(23)}.

%%K t12-06-1 p
6. --- Como ia adiantado, tive ocasião de ver o movimento dos viajantes
pela estação. Uns deixavam ou levantavam pequenos volumes no
\VUgras{depósito} \textsuperscript{(25)}; outros consultavam os \VUgras{horários}
\textsuperscript{(26)}. Os viajantes que chegavam apressavam-se para a
\VUgras{saída} \textsuperscript{(32)} com a sua \VUgras{mala} \textsuperscript{(24)} na
mão. Alguns esperavam na \VUgras{entrega das bagagens} \textsuperscript{(27)} e
apresentavam o seu \VUgras{talão} \textsuperscript{(29)} para retirar as suas malas,
\VUgras{caixas} \textsuperscript{(28)} ou \VUgras{cestos} \textsuperscript{(30)}.

%%K t12-07-1 p
7. --- Os \VUgras{empregados do fisco} \textsuperscript{(33)} revistavam com cuidado
os volumes para ver se havia algo a declarar.

%%K t12-08-1 p
8. --- Alguns viajantes dirigiam-se ao \VUgras{restaurante}
\textsuperscript{(34)}, onde um \VUgras{criado} \textsuperscript{(35)} se apressava a
servi-los. Tinham de se despachar, porque o \VUgras{comboio}
\textsuperscript{(36)}, que era um expresso, não ficava muito tempo na estação.

%%K t12-09-1 p
9. --- Passei depois ao cais de partida e procurei uma \VUgras{carruagem}
\textsuperscript{(37)} direta para não ter de mudar durante a viagem. Subi para
um carro de corredor que tem um \VUgras{compartimento} \textsuperscript{(38)} para
fumadores e outro reservado a «senhoras sós».

%%K t12-tit-2 sub
\VUpk{10.2pt}{40}{Partida de noite.}
\VUsaut{3.0mm}

%%K t12-10-1 p
10. --- Depois de subir o \VUgras{estribo} \textsuperscript{(40)}, fechar a
\VUgras{portinhola} \textsuperscript{(39)}, levantar a \VUgras{cortina}
\textsuperscript{(41)} e colocar os meus volumes pequenos na \VUgras{rede}
\textsuperscript{(43)}, sentei-me no \VUgras{banco} \textsuperscript{(42)}. Ao ver o
\VUgras{lampadeiro} \textsuperscript{(44)} acender os candeeiros, pensei que tinha
uma noite a passar na carruagem e chamei o empregado encarregado do
aluguer das \VUgras{almofadas} \textsuperscript{(45)} para lhe pedir uma.

%%K t12-11-1 p
11. --- O condutor do comboio veio verificar os bilhetes. Perguntei-lhe
por que não partíamos, visto que era a hora. Respondeu-me que o
\VUgras{chefe do comboio} \textsuperscript{(46)} estava ocupado a mandar colocar
umas bicicletas no \VUgras{furgão} \textsuperscript{(47)}. Além disso, ainda não se
tinham mudado todos os \VUgras{aquecedores de pés} \textsuperscript{(48)}.

%%K t12-12-1 p
12. --- Mas não íamos tardar a partir, porque o \VUgras{fogueiro}
\textsuperscript{(50)}, sobre o seu \VUgras{tênder} \textsuperscript{(49)}, pegava em
carvão para avivar o fogo da \VUgras{locomotiva} \textsuperscript{(54)}, e \VUgras{o
maquinista} \textsuperscript{(51)} só esperava o sinal do \VUgras{chefe adjunto de
estação} \textsuperscript{(52)}, que estava no \VUgras{cais} \textsuperscript{(53)}.

%%K t12-13-1 p
13. --- A \VUgras{caldeira} \textsuperscript{(56)} da locomotiva rugia surdamente; a
\VUgras{chaminé} \textsuperscript{(55)} vomitava torrentes de fumo misturado com
\VUgras{vapor} \textsuperscript{(60)}. Soou um \VUgras{apito} \textsuperscript{(59)} agudo e
os \VUgras{êmbolos} \textsuperscript{(58)} puseram-se em movimento, acionando as
\VUgras{rodas} \textsuperscript{(57)} da pesada máquina.

%%K t12-tit-3 sub
\VUsaut{3.0mm}
\VUpk{10.2pt}{40}{A caminho!}
\VUsaut{3.0mm}

%%K t12-14-1 p
14. --- A velocidade do comboio, que corria sobre os \VUgras{carris}
\textsuperscript{(64)}, acelerou-se; os \VUgras{cais} \textsuperscript{(65)} acabaram-se;
passámos as \VUgras{retretes} \textsuperscript{(66)} e também uma fila de carruagens
postas de parte noutra \VUgras{linha} \textsuperscript{(63)}: \VUgras{carruagens-cama}
\textsuperscript{(67)}, \VUgras{carruagem-restaurante} \textsuperscript{(68)},
\VUgras{carruagem-correio} \textsuperscript{(69)}.

%%K t12-noto-1 noto
\VUnotes{143.2mm}{%
\textit{(*) Nota.} --- \textit{Entende-se que a viagem se descreve tal
como era ao longo da fronteira de antes da guerra.}%
}

%%K t12-15-1 p
15. --- Depois passou diante dos nossos olhos a \VUgras{estação de
mercadorias} \textsuperscript{(71)}. Sob os alpendres, uns empregados carregavam
as \VUgras{mercadorias} \textsuperscript{(72)} enviadas em pequena velocidade. Uns
\VUgras{moços de manobras} \textsuperscript{(76)}, junto ao \VUgras{depósito de
água} \textsuperscript{(73)}, manobravam \VUgras{vagões de gado} \textsuperscript{(75)} e
\VUgras{vagões plataforma} \textsuperscript{(79)} para formar um \VUgras{comboio de
mercadorias} \textsuperscript{(80)}.

%%K t12-16-1 p
16. --- Passámos a última \VUgras{agulha} \textsuperscript{(78)}, junto à qual
estava o agulheiro; deixámos atrás o \VUgras{disco de sinais}
\textsuperscript{(86)} e a estação perdeu-se ao longe.

%%K t12-17-1 p
17. --- Em breve entrámos na \VUgras{ponte de ferro} \textsuperscript{(74)} que
atravessa o \VUgras{rio} \textsuperscript{(81)}. Passada a ponte, seguimos o curso
do rio, pelo qual potentes \VUgras{rebocadores} \textsuperscript{(82)} arrastavam
pesadas \VUgras{barcaças} \textsuperscript{(83)}. Vimos ainda um \VUgras{barco de
passageiros} \textsuperscript{(84)} no momento em que ia atracar ao
\VUgras{pontão} \textsuperscript{(85)}.

%%K t12-18-1 p
18. --- Depois de passar ao largo de várias estações, atravessámos alguns
\VUgras{túneis} \textsuperscript{(87)} e deixámos atrás de nós inumeráveis
\VUgras{casotas} \textsuperscript{(88)} de \VUgras{guarda-barreiras}. Embalado pelo
solavanco do comboio, adormeci profundamente.

%%K t12-tit-4 sub
\VUsaut{3.0mm}
\VUpk{10.2pt}{40}{A estação de Deutsch-Avricourt.}
\VUsaut{3.0mm}

%%K t12-19-1 p
19. --- Acordou-me de repente a voz de um empregado que gritava:
«Deutsch-Avricourt! \textsuperscript{(*)} Desçam todos!». Era a inspeção
aduaneira e todas as fastidiosas formalidades da fronteira. Tive de pôr o
meu relógio a horas pelo \VUgras{relógio} \textsuperscript{(89)} da estação, que ia
cinquenta e cinco minutos adiantado; mal acordado, distinguia com
dificuldade o \VUgras{ponteiro grande} \textsuperscript{(90)} do \VUgras{pequeno}
\textsuperscript{(91)}; o próprio \VUgras{mostrador} \textsuperscript{(92)} estava de todo
confuso pelo nevoeiro da manhã.

%%K t12-20-1 p
20. --- Enquanto os aduaneiros faziam a sua inspeção, decifrei a bocejar
os \VUgras{cartazes} \textsuperscript{(93)} que decoram as paredes da estação.
Tomei o pequeno-almoço para matar o tempo, consultei o mapa da
\VUgras{rede} \textsuperscript{(94)} alemã para ver que distância me separava ainda
de Estrasburgo, e voltei a subir para a carruagem, reconfortado pela
esperança de chegar em breve ao termo da minha viagem.
\VUsaut{6.0mm}
\VUfilet{22mm}
